\section{	exorpdfstring{THE ``DICTIONARY'' LIE GROUPS --- LIE ALGEBRAS.}{THE “DICTIONARY” LIE GROUPS — LIE ALGEBRAS.}}

The theory of ``finite and continuous'' groups, developed by Lie in numerous memoirs starting in 1874, is set out systematically in the imposing treatise ``Theorie der Transformationsgruppen'' ({[}203{]}, 1888-1893), written in collaboration with F. Engel; \(^{7}\) it is the object of the first volume and of the last five chapters of the third, the second being consecrated to contact transformations.

As the title indicates, there was never any question in this work of anything other than transformation groups, in the meaning of the equations (25.4), where the space of ``variables'' \(x_{i}\) and the space of ``parameters'' \(a_{j}\) initially both play equally important roles. Besides the concept of ``abstract'' group is not clearly expressed at this time; when in 1883 ({[}202{]}, vol.~5, Abh. XII, pp.~311-313) Lie remarks that with the notation of (25.10), the equation \(w = H(u, v)\) , which gives the parameters of the product of two transformations of the group, defines a new group, it is as a transformation group on the space of parameters that he is considering it, obtaining thus what he calls the ``group of parameters'' (he even obtains two of them, which are none other than the group of left translations and the group of right translations). \(^{8}\)

The variables \(x_{i}\) and the parameters \(a_{j}\) in the equations (25.4) are assumed complex in principle (except in chaps. XIX-XXIV of v. 3), and the functions \(f_{i}\) analytic; Lie and Engel are of course conscious of the fact that these functions are not in general defined for all complex values of the \(x_{i}\) and the \(a_{j}\) and that, in consequence, the composition of such transformations raises serious difficulties ({[}203{]}, v. I, pp.~15-17, pp.~33-40 and passim); and although, later, they express themselves almost always as if composition of the transformations that they study were always possible without restriction, it is no doubt only for the convenience of the statements, and they re-establish explicitly the ``local'' point of view each time that it is necessary (cf.~loc. cit., pp.~168 or 189 for example or ibid., v. 3, p.~2, note at the bottom of the page); in other words, the mathematical object that they are studying is near to what we would now call a partial operation law. They do not fail to consider, on occasion, global groups, for example the 4 series of classical groups ({[}203{]}, v. 3, p.~682), but do not seem to have set themselves the question of what in general a ``global group'' can be; it suffices for them to be able to obtain, for the ``parameters'' of the classical groups (the ``variables'' of these groups do not introduce any difficulty, since it is a case of linear transformations of \(\mathbf{C}^n\) ), systems of ``local'' parameters in the neighbourhood of the identical transformation, without their being concerned about the domain of validity of the formulae that they write down. They set themselves however one problem that stands out clearly from the local theory:⁹ the study of ``mixed'' groups, that is to say groups having a finite number of connected components, such as the orthogonal group ({[}203{]}, v. I, p.~7). They present this study as that of a set of transformations closed under composition and taking inverses, which is the union of sets \(H_{j}\) such that each is described by systems of functions \(f_{i}^{(j)}\) as in (25.4); the number of (essential) parameters of each \(H_{j}\) is even assumed a priori to depend on j, but they show in fact that this number is the same for all the \(H_{j}\) . Their main result is then the existence of a finite continuous group G such that \(H_{j} = G \circ h_{j}\) for an \(h_{j} \in H_{j}\) and for all j; they establish also that G is normal in the mixed group and remark that the determination of the invariants of this latter reduces to that of the invariants of G and of a discontinuous group ({[}203{]}, v. I, chap.~18).

The general theory developed in {[}203{]} finishes up (without that being stated in a very systematic way by the authors) by forging a ``dictionary'' making the translation from properties of ``finite continuous'' groups to those of the set of their infinitesimal transformations. It is based on the ``three theorems of Lie'', each one of which is made up of an assertion and its converse.

The first theorem ({[}203{]}, v. I, pp.~33 and 72 and v. 3, p.~563) affirms in the first instance that if in (25.4) the parameters are effective, the functions \(f_{i}\) satisfy a system of partial differential equations of the form

\[
\frac {\partial f _ {i}}{\partial a _ {j}} = \sum_ {k = 1} ^ {r} \xi_ {k i} (f (x, a)) \psi_ {k j} (a) (1 \leq i \leq n)\tag{25.15}
\]

where the matrix \((\xi_{kl})\) is of maximum rank and \(\det(\psi_{kj}) \neq 0\) ; reciprocally, if the functions \(f_{i}\) have this property, the formulae (25.4) define a groupuscule of transformations.

The second theorem ({[}203{]}, v. I, pp.~149 and 158, and v. 3, p.~590) gives relations between the \(\xi_{ki}\) on the one hand, the \(\psi_{ij}\) on the other: the conditions on the \(\xi_{ki}\) can be written in the form

\[
\sum_ {k = 1} ^ {n} \left(\xi_ {i k} \frac {\partial \xi_ {j l}}{\partial x _ {k}} - \xi_ {j k} \frac {\partial \xi_ {i l}}{\partial x _ {k}}\right) = \sum_ {k = 1} ^ {r} c _ {i j} ^ {k} \xi_ {k l} (1 \leq i, j \leq r, 1 \leq l \leq n)\tag{25.16}
\]

where the \(c_{ij}^{k}\) are constants \((1 \leq i, j, k \leq r)\) anti-symmetric in i, j. The conditions on the \(\psi_{ij}\) in the form given by Maurer, are:

\[
\frac {\partial \psi_ {k l}}{\partial a _ {m}} - \frac {\partial \psi_ {k m}}{\partial a _ {l}} = \frac {1}{2} \sum_ {1 \leq i, j \leq r} c _ {i j} ^ {k} (\psi_ {i l} \psi_ {j m} - \psi_ {j l} \psi_ {i m}) (1 \leq k, l, m \leq r).\tag{25.17}
\]

By introducing the matrix \((\alpha_{ij})\) contragredient to \((\psi_{ij})\) and the infinitesimal transformations

\[
X _ {k} = \sum_ {i = 1} ^ {n} \xi_ {k i} \frac {\partial}{\partial x _ {i}}, A _ {k} = \sum_ {j = 1} ^ {r} \alpha_ {k j} \frac {\partial}{\partial a _ {j}} (1 \leq k \leq r),\tag{25.18}
\]

(25.16) and (25.17) can be written respectively:

\[
[ X _ {i}, X _ {j} ] = \sum_ {k = 1} ^ {r} c _ {i j} ^ {k} X _ {k} (1 \leq i, j \leq r)\tag{25.19}
\]

\[
\left[ A _ {i}, A _ {j} \right] = \sum_ {k = 1} ^ {r} c _ {i j} ^ {k} A _ {k}.\tag{25.20}
\]

Conversely, if given r infinitesimal transformations \(X_{k}(1 \leq k \leq r)\) , linearly independent and satisfying conditions (25.19), the one parameter subgroups generated by these transformations generate a group of transformations with r essential parameters.

Finally, the third theorem ({[}203{]}, v. I, pp.~170 and 297 and v. 3, p.~597) reduces the determination of systems of infinitesimal transformations \((X_{k})_{1\leq k\leq r}\) satisfying (25.19) to a purely algebraic problem: we must have

\[
c _ {i j} ^ {k} + c _ {j i} ^ {k} = 0\tag{25.21}
\]

\[
\sum_ {l = 1} ^ {r} (c _ {i l} ^ {m} c _ {j k} ^ {l} + c _ {k l} ^ {m} c _ {i j} ^ {l} + c _ {j l} ^ {m} c _ {k i} ^ {l}) = 0 (1 \leq i, j, k, m \leq r).\tag{25.22}
\]

Conversely, \(^{10}\) if (25.21) and (25.22) are satisfied, there exists a system of infinitesimal transformations satisfying relations (25.19), from whence a group of transformations with r parameters (in other words, the linear combinations with constant coefficients of the \(X_{k}\) make up a Lie algebra, and conversely every Lie algebra of finite dimension can be obtained in this way).

These results are completed by the study of isomorphism questions. Two groups of transformations are said to be similar if it is possible to pass from one to the other by an invertible transformation of co-ordinates on the variables and an invertible transformation of co-ordinates on the parameters: from the beginning of his research, Lie had met this notion in a natural way in connection with the definition of the ``canonical parameters''. He shows that two groups are similar if, by means of a transformation on the ``variables'', one can bring the infinitesimal transformations of the one onto those of the other ({[}203{]}, v. I, p.~329). A necessary condition for this to be so is that the Lie algebras of the two groups be isomorphic, which Lie expresses by saying that the groups are ``gleichzusammengesetzt''; but this condition is not sufficient, and a whole chapter ({[}203{]}, v. I, chap.~19) is devoted to obtaining supplementary conditions insuring that the groups should be ``similar''. On the other hand the theory of permutation groups supplied the notion of the ``holoedric isomorphism'' of two such groups (isomorphism of the underlying ``abstract'' groups); Lie transposes this notion to groups of transformations, and shows that two such groups are ``holoedric isomorphic'' if and only if their Lie algebras are isomorphic ({[}203{]}, v. I, p.~418). In particular, every group of transformations is holoedrically isomorphic to each of its groups of parameters, and that shows that, when one wishes to study the structure of the group, the ``variables'' on which they operate do not matter much and that in fact everything reduces to the Lie algebra. \(^{11}\)

Always by analogy with the theory of permutation groups. Lie introduces the notions of subgroups, normal subgroups, ``meriedric isomorphisms'' (surjective homomorphisms), and shows that they correspond to that of subalgebras, ideals and surjective homomorphisms of Lie algebras; besides he had met very early on a particularly important example of ``meriedric isomorphism'', the adjoint representation, and had recognised its links with the centre of the group ({[}202{]}, vol.~5, Abh. III, pp.~42-75). For these results, as for the fundamental theorems, the essential tool is the theorem of Jacobi-Clebsch giving the complete integrability of a differential system (one of the forms of the theorem called ``of Frobenius''); he gives a new proof of it, as it happens, using groups with one parameter ({[}203{]}, v. I, chap.~6).

The notions of transitivity and of primitivity, so important for groups of permutations, arise also naturally for the ``finite and continuous'' groups of transformations, and the treatise of Lie-Engel makes a detailed study of them ({[}203{]}, v. I, chap.~13 and passim); the relations with the subgroups stabilising a point and the notion of homogeneous space are observed (for as much as they could be without putting oneself in a global point of view) ({[}203{]}, v. I, p.~425).

Finally, the ``dictionary'' is completed, in {[}203{]}, by the introduction of the notions of derived group and soluble group (called ``integrable group'' by Lie; this terminology, suggested by the theory of differential equations, will remain in use until the work of H. Weyl) ({[}203{]}, v. I, p.~261 and v. 3, pp.~678-679); the relation between commutators and brackets had besides been observed by Lie already in 1883 ({[}202{]}, vol.~5, p.~358).

